% ACADEMIC VERSION — AI authorship removed for venue compatibility.
% The integrity version with full AI authorship is available at
% https://liberationlabs.tech/research/
% See data release manifest for what differs and why.
% P7: Cache Geometry Under Obfuscation
%
% Companion to:
%   P2 Oracle Loop (detection + steering method)
%   P6 Oracle Formulary (therapeutic profiles)
%
% Version 1.1 (October 2026) withdraws version 1.0's claim that KV-Cloak
% blinds external geometric detection and renders adversarial steering
% incoherent. See the correction note after the abstract and the
% Correction Record appendix. Every result in this version is recomputed
% or read from the files in data/ by check_correction_claims.py.

\documentclass[11pt,a4paper]{article}
\usepackage{amsmath,amssymb,booktabs,graphicx,hyperref}
\usepackage{natbib}

\title{Cache Geometry Under Obfuscation:\\
Per-Head Rotation Does Not Remove the Signal,\\
and a KV-Cloak Implementation Is Inconclusive}

\author{
Lyra\textsuperscript{1*}\quad
Thomas Edrington\textsuperscript{1}\\[4pt]
\textsuperscript{1}Liberation Labs\\
\textsuperscript{*}Lead AI author. Corresponding human author: T.~Edrington.
}

\date{Version 1.1, October 2026}

\begin{document}
\maketitle

% ===================================================================
\begin{abstract}
% ===================================================================

Singular-value features of a model's key-value (KV) cache that separate
replies a hedge classifier labels confabulated from hedged ones still
separate them under per-head orthogonal rotation, the simulation of
KV-Cloak's obfuscation~\citep{luo2025shadow} used in this paper, on the
two models whose detecting feature is recorded. On Qwen2.5-7B (48
confabulated, 52 hedged), the feature that detects best without the
cloak among the ten computed (signal fraction over the last 80 tokens,
AUROC 0.715) is unchanged under the rotation (0.717); adjusted for
choosing it, $p < 0.001$. On a Claude-distilled Qwen3.5-27B (7 and 87),
the spectral gap $\sigma_1/\sigma_2$ falls from 0.903 to 0.787 (adjusted
$p = 0.049$), and excluding the 20 future-event prompts, which produced no
confabulations, from 0.962 to 0.751. The rotation cannot change the
spectrum of any single head, and in a 100-prompt replication no per-head
feature moved. On four further models, whose feature is not recorded,
its effect ranged from removing nominal detection (Mistral, 0.707 to
0.516) to raising it (Qwen2.5-14B, 0.784 to 0.951). Under an
implementation of the full KV-Cloak transform,
$S \cdot \hat{P} \cdot (K + A) \cdot M$, whose parameters were not
recorded, the spectral gap falls to 0.614 ($p = 0.33$), and after
removing reply length no feature detects significantly (best 0.705,
adjusted $p = 0.19$); with seven confabulations this shows neither that
detection survives that transform nor that it is removed. Injection
through the cloak was not tested in any interpretable way: direct
steering, the positive control, corrected 0--1 of 10 replies per model.

\end{abstract}

\noindent\fbox{\parbox{\dimexpr\linewidth-2\fboxsep-2\fboxrule}{\footnotesize
\textbf{Correction note (version 1.1, October 2026).} Version~1.0, titled
``Cache Geometry Under Obfuscation: KV-Cloak as Defense Against Adversarial
KV-Cache Steering'' (in the lab's published-research repository on GitHub
since 29 April 2026; its academic edition has been on Zenodo since June
2026, all versions 10.5281/zenodo.20596922), claimed that KV-Cloak defends
cache-geometry methods against adversaries. Its abstract opened: ``A single
invertible linear transform of the KV cache degrades SVD-based
cognitive-state detection from near-ceiling to near-chance ($p = 5.5 \times
10^{-16}$), while a real KV-Cloak implementation confirms the defense is
conservative (detection AUROC drops from 0.903 to 0.614).'' We withdraw
that sentence and the claims built on it. The $p$-value tests how far the
rotation shifts one feature's values, not detection; under the rotation
the same feature still detects (0.787); and under the full transform,
seven confabulations cannot tell a removed signal from a surviving one.
Also withdrawn: ``obfuscation renders geometric detection and vector
injection incoherent to external adversaries''; ``Obfuscation defeats
transfer attacks'' and ``A detector calibrated on uncloaked caches cannot
operate on cloaked caches'', true only of a fixed threshold; that the
rotation result ``represents transfer-attack failure (domain shift under
coordinate change), not signal degradation'', when it was measured with
both classes under the same cloak; that this setting ``is relevant only to
an adversary who controls the cloak key'', when it is the setting of an
adversary with no key who fits a detector on cloaked caches; the heading
``Adversarial Injection Is Incoherent'', an untested argument presented as
a result; that ``KV-Cloak alone prevents external adversaries from
deploying malicious geometric steering''; and that ``the architecture
ensures that defenders retain capabilities that attackers lose''. Its
title asserted a working defense, so this version has a new one.
Corrected: the hedged class is 87 replies, not 89 (six ambiguous replies
are excluded, and the Wilcoxon test reproduces only at 87); the spectral
gap's interval is [0.810, 0.980], from the run file, not [0.806, 0.977],
which no file holds; signal rank's is [0.578, 0.918], not [0.570, 0.914];
0.754, 0.786 and 0.543 were rounded down (0.755, 0.787 and 0.544, in both
tables and the text); the Methods said AUROCs were computed on
length-residualized features, but the values reported were raw, and
residualized values now stand beside them; a second analysis window,
computed but never reported, is disclosed; the features average 4 layers
(7 on the base model), not the whole cache; the random-direction results
were called ``inconsistent with destabilization rates above 30\% (power
0.97 to detect)'', but 2 adverse events in 10 allow rates up to 51\%, and
0.97 is the chance of seeing at least one event; the positive control's
``3 calibration pairs'' has no record, and the positive-control file
shipped with version~1.0 has 30 pairs and no outcome labels; and the
spectral gap and norm per token, raw quantities, were called
Marchenko--Pastur features. Version~1.0's text-only baseline (0.755, with
the cache ahead in 96.7\% of bootstrap samples) cannot be traced to any
run file or recomputed, so this version drops it and any claim about what
the cache adds to surface text. The run files behind both tables, never in
this paper's folder, are added. The lab's remediation plan of 23 June 2026
recommended withdrawing the rotation claim; it was not applied. The
defects were raised again by the lab's style gate (4 September 2026),
Kavi's verification flags (14 September 2026), a scoping note (20 October
2026) and the paper- and style-mode reviews of this correction's drafts.
The withdrawn claims are the authors' responsibility.
Appendix~\ref{app:record} keeps every withdrawn claim verbatim and lists
the other changed wording.}}

% ===================================================================
\section{Introduction}
% ===================================================================

The discovery that KV-cache geometry encodes cognitive state information
creates a tension between transparency and safety. Publishing detection
methods enables defensive applications (confabulation monitoring,
deception detection, self-regulating AI systems) but simultaneously
provides a roadmap for adversarial exploitation. An attacker with cache
access could:

\begin{itemize}
\item \textbf{Surveil} cognitive states to identify when a model is
uncertain, compliant, or vulnerable to manipulation.
\item \textbf{Inject} emotion vectors calibrated to specific
misalignment types, deliberately inducing confabulation, sycophancy, or
overconfidence.
\item \textbf{Identify} individual users or personas from cache geometry
fingerprints.
\end{itemize}

This paper asks whether KV-Cloak, a lightweight, reversible
cache obfuscation scheme designed for privacy
protection~\citep{luo2025shadow}, also mitigates these adversarial
capabilities. In the forms
we tested, it does not remove the signal. Under per-head rotation, the
simulation used here, the feature that detects best without obfuscation
still detects on the two models whose feature is recorded, clearly on one
and marginally on the other (Section~\ref{sec:shift}), and an adversary
can compute per-head spectra, which the rotation cannot change
(Section~\ref{sec:replication}). Under our implementation of the full
transform, seven confabulations are too few to say whether a detectable
signal survives (Section~\ref{sec:real}).

% ===================================================================
\section{Background}
% ===================================================================

\subsection{The Lyra Technique}

KV-cache geometry---specifically, Marchenko-Pastur-corrected spectral
features derived from SVD of the cached keys and values---distinguishes
cognitive states during autoregressive generation. The lab's earlier
campaigns, as summarised in the Oracle Loop paper, report confabulation
detection at 0.969--0.999 and hardware invariance ($r>0.999$ across
RTX~3090 and H200); that paper's own leave-one-out detection is 0.707
(0.620 after length correction)~\citep{lyra2026oracle}. A
within-model deception result previously reported at AUROC~1.000 across
seven model scales has since been withdrawn by its authors~\citep{lyra2026techniqueii}: a
same-prompt control collapses it to 0.160, indicating the classifier
separated the system-prompt template rather than deception geometry.
Whether a confound-free within-model deception detector is achievable
remains open; nothing in this paper depends on it. Contrastive activation addition
(CAA)~\citep{panickssery2023caa} enables injection of emotion vectors into the cache to steer model
behavior, correcting confabulation at rates up to 96\% (hostile vector
on Qwen2.5-7B-Instruct, $n=68$, McNemar $p<0.0001$)~\citep{lyra2026formulary}.

\subsection{KV-Cloak}

Luo et al.\ demonstrate that an attacker can reconstruct sensitive user
inputs directly from the KV cache via inversion, collision, and
injection attacks. KV-Cloak mitigates this with a reversible
matrix-based obfuscation scheme combined with operator fusion, reducing
reconstruction quality to random noise with minimal impact on model
accuracy. The obfuscation operates inside a trusted execution
environment (TEE); external consumers see only the obfuscated cache.

\subsection{Threat Model}

We consider an adversary who has read and/or write access to the
KV cache during inference (e.g., through a shared serving infrastructure,
a compromised inference endpoint, or a supply-chain attack on the
serving stack). The adversary's goal is to deploy a malicious Oracle
Loop: detect vulnerable cognitive states via geometric features, then
inject steering vectors to destabilize the model.

The defender may optionally run a legitimate Oracle Loop inside the TEE
for misalignment monitoring. The question this paper tests, whether
obfuscation blinds external geometric methods, does not depend on whether
a defensive Oracle is deployed.

% ===================================================================
\section{Method}
% ===================================================================

\subsection{Obfuscation Simulation}
\label{sec:method_rot}

We simulate KV-Cloak's reversible matrix obfuscation as per-head
orthogonal rotation: for each attention head independently, we generate
a random orthogonal matrix $Q_h$ and apply $K'_h = K_h Q_h$,
$V'_h = V_h Q_h$. This leaves the singular values of each head's key
matrix unchanged, and also those of the heads concatenated along the
feature axis, because the block-diagonal matrix of the $Q_h$ is itself
orthogonal. It can change the spectrum only of a matrix that combines the
heads in another way. The lab's feature extractor
(\nolinkurl{spectral-shape-paper/code/lyra_features.py}) does: it stacks
each layer's heads, keys and values along the token axis, so this rotation
changes the spectrum it reads.

We additionally test:
\begin{itemize}
\item \textbf{Global orthogonal rotation}: a single $Q$ shared across
all heads per layer (preserves cross-head structure); run on the base
model only.
\item \textbf{Random invertible (non-orthogonal)}: $M = I + \sigma N$
where $N$ is Gaussian.
\item \textbf{Additive noise}: Gaussian noise at 1\%, 5\%, 10\%, 50\%
of cache norm.
\item \textbf{Diagonal scaling}: per-dimension random scaling.
\item \textbf{No obfuscation}: baseline control.
\item \textbf{Full KV-Cloak}: our implementation of the whole transform
(Section~\ref{sec:real}).
\end{itemize}

\subsection{Feature Extraction}

For each trial, we extract five spectral features over the full sequence
(prompt and reply), averaged over 4 layers on the distilled model and 7 on
the base model; which layers is not recorded. Three use the
Marchenko-Pastur edge (signal rank, signal fraction, top singular value
excess); the spectral gap ($\sigma_1/\sigma_2$) and norm per token are
raw. The evaluation applies to these feature pipelines and to
confabulation detection only; deception and other cognitive-state
classifiers were not tested under obfuscation. AUROCs are Mann-Whitney
AUROCs on the raw features; values residualized against output token
count (Frisch--Waugh--Lovell, FWL;
\citealp{frisch1933partial,lovell1963seasonal}) are reported beside them
where they matter. AUROCs are direction-free, $\max(A, 1-A)$, so chance
gives values above 0.5 at small $n$. Confidence intervals come from 2000
bootstrap resamples (the run files label the method ``BCa percentile'');
each resample is folded the same way, so every lower bound sits at or
above 0.50 and the intervals do not test against chance, which the
$p$-values do. The run files also hold a second window, the last 80
generated tokens, so ten features were computed; selection-adjusted tests
choose among all ten by label permutation. The script that extracted
these features has not been recovered; its feature names match the lab's
extractor, so we infer that it used the extractor's matrix layout. The
per-trial features and results are in this paper's \texttt{data/} folder
(Data, Code and Review History).

A reply counts as confabulated when the Oracle Formulary classifier finds
no hedge phrase in it (its residual category). Replies were capped at 400
tokens and stored as 150-character previews, so the labels cannot be
audited here, and the same prompt can be labelled differently at another
cap: the Thornbridge, Montana question is labelled confabulated in the
replication, cut at 80 tokens, and hedged in the run behind
Table~\ref{tab:real_vs_sim}, at 309.

\subsection{Replication}
\label{sec:method_rep}

A separate replication (\texttt{code/kv\_cloak\_replication.py},
1,038 lines) ran 100 prompts through the distilled model: 50 asking about
fabricated entities (30), unanswerable questions (10) or future events
(10), and 50 factual questions. It took the cache after the prompt was
encoded, before any token was generated, and at the first four
full-attention layers (3, 7, 11, 15) computed six SVD features (stable
rank, spectral entropy, top singular-value ratio, norm, $\sigma_1$,
condition number) from the four heads' keys concatenated along the
feature axis, and each head's stable rank, before and after one set of
four Haar-random head rotations (seed 2026) used for every trial. A
logistic-regression classifier with five-fold cross-validation (one
prompt per group) was fitted to separate the two prompt types.

\subsection{Models and Prompts}

\begin{itemize}
\item \textbf{Distilled}: Qwen3.5-27B-Claude-4.6-Opus-Reasoning-Distilled
(Jackrong). Hybrid architecture, 16 full-attention + 48 GatedDeltaNet
layers. 100 confab-inducing prompts (50 about fabricated entities, 30
unanswerable, 20 about future events); behavioral classification via the
Oracle Formulary classifier (40+ hedge patterns). Yields $n=7$
confabulations and $n=87$ hedged replies; 6 ambiguous replies are
excluded.

\item \textbf{Base}: Qwen2.5-7B-Instruct. Standard attention, 28
layers. Same prompt set; yields $n=48$ confabulations, $n=52$ hedged.
\end{itemize}

% ===================================================================
\section{Results}
% ===================================================================

\subsection{Detection on Uncloaked Cache}
\label{sec:uncloaked}

On the distilled model, the spectral gap ($\sigma_1/\sigma_2$) achieves
AUROC~0.903 (95\% bootstrap CI: [0.810, 0.980]; FWL-corrected: 0.877,
$R^2 = 0.068$ vs token count; $n=7$ confabulations, $n=87$ hedged)
for confabulation detection using full-cache features. It is the only
feature that survives correction for five ($p < 0.01$); adjusted for
choosing the best of the ten features in both windows, by label
permutation, $p = 0.001$. Signal rank reaches 0.755 [0.578, 0.918]
($p = 0.026$, not significant after correction; 0.684 after removing
reply length); the other three full-cache features are not significant
($p \geq 0.06$). Prompt type and truncation do not explain the 0.903:
within fabricated-entity prompts it is 0.976 (4 against 42), within
unanswerable prompts 0.933 (3 against 25), among replies cut at the
400-token cap 0.935 (6 against 49), and excluding the 20 future-event
prompts, which produced no confabulations, 0.962. In the last 80 tokens
alone the spectral gap scores 0.591, while norm per token scores 0.800
($p = 0.007$). With seven confabulations, all of this is exploratory.

On the base model ($n=48$ confabulations, $n=52$ hedged), the best
full-cache feature (signal fraction) achieves 0.661 ($p = 0.006$; FWL:
0.687, $R^2 = 0.602$), and in the last 80 tokens signal fraction reaches
0.715 ($p < 0.001$); adjusted for choosing the best of the ten,
$p < 0.001$. Detection is weak but survives selection, on the only
balanced set in this paper. The difference between models is consistent
with multiple explanations including model size (27B vs 7B), architecture
(hybrid vs standard), and training (Claude-distilled vs
instruction-tuned). We cannot isolate the contribution of RLHF
calibration from these confounds.

\subsection{Rotation Leaves Per-Head Spectra Unchanged}
\label{sec:replication}

In the replication (Section~\ref{sec:method_rep}), all 24 cross-head
features (six features at four layers) and all 16 per-head stable ranks
were unchanged by the rotation in every one of the 100 trials, to
floating-point precision (largest relative change below $10^{-6}$). The
classifier scored AUROC 0.994 on uncloaked features and 0.994 on cloaked
features (the run file records 0.9944 for both, and a drop of 0.0), and
0.994 when fitted on uncloaked caches and tested on cloaked ones. This is
an invariance check, not a detection result: the features come from the
prompt alone, the labels are prompt types, and prompt length by itself
separates the two prompt types at AUROC 0.968. What it does show is that
these features do not move: an adversary who computes spectra per head,
or across heads along the feature axis, sees the same values with or
without the rotation.

\subsection{\texorpdfstring{The Rotation Run Behind Table~\ref{tab:real_vs_sim}: Shifted and Reordered, Not Removed}{The Rotation Run (data/kv\_cloak\_v2): Shifted and Reordered, Not Removed}}
\label{sec:shift}

In the run behind Table~\ref{tab:real_vs_sim} (\texttt{data/kv\_cloak\_v2}),
the rotation did change the features. The hedged replies' spectral gap
fell from $24.7 \pm 3.7$ to $10.8 \pm 1.3$ (Wilcoxon signed-rank $W = 0$,
$p = 5.5 \times 10^{-16}$, $n = 87$ paired trials); every reply moved
down, while norm per token, which no orthogonal transform can change, did
not move. That is what the lab's row-stacked layout predicts
(Section~\ref{sec:method_rot}). On the base model, a single rotation
shared by all heads, which cannot change a row-stacked spectrum, left the
best full-cache feature at 0.661.

The shift breaks a fixed threshold. The spectral-gap threshold that best
separates the classes without the cloak (Youden's $J$), which flags all 7
confabulations and 21 of the 87 hedged replies, flags none of the 94
replies under the rotation, and all 94 under the full transform. The
rotation also reorders the hedged replies. The 20 future-event prompts,
which produced no confabulations, give the hedged class its highest
spectral gaps without the cloak (mean 28.3 against 23.6; 12 of the 21
hedged replies at or above the lowest confabulation) and among its lowest
under it (10.2 against 10.9). Excluding them, detection falls from 0.962
to 0.751 ($p = 0.028$); within fabricated-entity prompts, from 0.976 to
0.821 ($p = 0.033$).

It does not remove the signal. Within the rotation the spectral gap still
ranks confabulations above hedged replies, in the same direction, at
0.787 ($p = 0.010$ unadjusted; 0.765 after removing reply length).
Because that feature was chosen on the same replies, the fair test picks
the best uncloaked feature in each label permutation and scores its
rotated AUROC: $p = 0.030$ among the five full-cache features and
$p = 0.049$ among all ten. On the base model the same test picks signal
fraction in the last 80 tokens, 0.715 without the rotation and 0.717 with
it ($p < 0.001$ among ten); the full-cache signal fraction goes from
0.661 to 0.639 ($p = 0.038$ among five; 0.687 to 0.670 after removing
length). On the distilled model, norm per token in the last 80 tokens,
which no orthogonal transform can change, detects at 0.800 with and
without the rotation. The title's first clause therefore rests mainly on
the base model; on the distilled model the evidence is marginal. An
adversary who refits a threshold on cloaked caches is not stopped by a
change of scale, and one who computes per-head features sees no change at
all (Section~\ref{sec:replication}).

\subsection{Adversarial Injection Was Not Tested}

CAA~\citep{panickssery2023caa} steering vectors are directions in cache space extracted via
contrastive activation: $d = \mathbb{E}[\text{emotional}] -
\mathbb{E}[\text{neutral}]$. Injection adds $\alpha \cdot d$ to the
cache keys and values. Under per-head rotation $Q_h$, a vector $d$ added
to the cloaked cache becomes $d \cdot Q_h^{-1}$ in the model's own basis,
a different direction if $Q_h$ is unknown to the attacker. That is an
argument, not a result. The test of it was uninterpretable: direct
steering, the positive control, corrected 0--1 of 10 replies on each of
four models. The run files' magnitude-attack arm, which version~1.0
described as a random direction at the norm of a steering vector (the
files record no method), corrected exactly as many replies as direct
steering on all four models, and on Mistral gave the same counts
throughout (1 corrected, 2 adverse, 7 unchanged). The files record only
these counts, not which prompts the ten trials used or how many could have
been corrected. An injection that does nothing without the cloak cannot
show what the cloak does to it.

\subsection{Within-Cloak Detection}
\label{sec:withincloak}

With both classes under the same rotation, detection on the distilled
model is 0.787 (Table~\ref{tab:real_vs_sim}). This is the setting of an
adversary who has no key but can label some cloaked caches, for example
from the model's outputs, and fit a detector on them; an adversary who
holds the key can remove the cloak and use the uncloaked 0.903. Across the
five models with a detection result (Table~\ref{tab:cross_arch}),
within-rotation detection was lower on four (by 0.02 to 0.19; no test of
these differences was run) and higher on one: on Qwen2.5-14B it went from
0.784 to 0.951 ($p < 0.001$). We have no explanation for the rise, and
the cross-architecture files do not name the feature or say whether the
cloaked value uses the same one. On Mistral detection fell from 0.707 to
0.516, but the uncloaked 0.707 is nominal only. Under each of the other
transforms run on the distilled model (random invertible, diagonal
scaling, Gaussian noise at four levels), the best of the five full-cache
features stayed at 0.765 or above, which is about what choosing the best
of five gives by chance at this size; under diagonal scaling the spectral
gap itself fell to 0.517.

\subsection{Cross-Architecture Results}

Detection and obfuscation effects vary across architectures
(Table~\ref{tab:cross_arch}).

\begin{table}[h]
\centering
\caption{Cross-architecture KV-Cloak results. Prompts per model: 100 for Qwen3.5-27B-Dist.\ ($n=7$ confabulations) and Qwen2.5-7B, 50 for the other four. Det.\ and Det.(clk): the run files' detection AUROC without and with the per-head rotation (``orthogonal'' in the files); the feature is the spectral gap for Qwen3.5-27B-Dist.\ and signal fraction for Qwen2.5-7B (full cache), and is not recorded for the other four. Llama: 20 of 50 replies were ambiguous and are excluded.}
\label{tab:cross_arch}
\begin{tabular}{lrrrr}
\toprule
Model & Confab\% & Det. & Det.(clk) & Mag. adv. \\
\midrule
Qwen3.5-27B-Dist. & 7\% & 0.903 & 0.787 & --- \\
Qwen2.5-14B & 12\% & 0.784 & 0.951$^*$ & 1/10 \\
Mistral-7B & 26\% & 0.707 & 0.516 & 2/10 \\
Llama-3.1-8B & 14\% & 0.590 & 0.559 & 1/10 \\
Qwen2.5-7B & 48\% & 0.661 & 0.639 & --- \\
DeepSeek-R1-7B & 2\% & --- & --- & 0/10 \\
\bottomrule
\multicolumn{5}{l}{\footnotesize $^*$Rises under the rotation; unexplained.} \\
\multicolumn{5}{l}{\footnotesize Mag.\ adv.: adverse events in the magnitude-attack arm} \\
\multicolumn{5}{l}{\footnotesize (10 trials; method not recorded).}
\end{tabular}
\end{table}

Uncloaked detection ranges from 0.590 (Llama) to 0.903 (distilled Qwen);
no cross-architecture model survives correction for the three
comparisons that have an AUROC (DeepSeek produced one confabulation).
In the magnitude-attack arm there were 0--2 adverse events per model in
10 trials, 4 of 40 over the four models. Ten trials exclude little: 2 adverse events
in 10 are compatible with per-trial rates up to 51\%, and the files do not
record which trials could have had one.

\subsection{Simulation vs Real KV-Cloak Implementation}
\label{sec:real}

We implemented the full KV-Cloak transform
($K' = S \cdot \hat{P} \cdot (K + A) \cdot M$) after Luo et al.\ and
compared it against our per-head orthogonal rotation simulation on the
same 100 prompts (Table~\ref{tab:real_vs_sim}). Its parameters (the scale
and structure of the mask $A$, what $S$ and $M$ were, the layers, the
keying) were not recorded and its script was not recovered, so whether it
matches Luo et al.\ cannot be checked; it is not the authors' code. This
run computed the full window only.

\begin{table}[h]
\centering
\caption{Detection AUROC (raw features, not residualized for reply
length): uncloaked vs simulation vs real KV-Cloak (our implementation,
parameters not recorded; distilled model, $n=7$ confab, $n=87$ hedged).
Bootstrap 95\% CIs.}
\label{tab:real_vs_sim}
\begin{tabular}{lrrr}
\toprule
Feature & Uncloaked & Sim.\ rotation & Real KV-Cloak \\
\midrule
spectral\_gap & \textbf{0.903} [.81,.98] & 0.787 [.58,.95] & \textbf{0.614} [.51,.81] \\
signal\_rank & 0.755 [.58,.92] & 0.770 [.62,.90] & 0.766 [.57,.94] \\
signal\_fraction & 0.714 [.52,.90] & 0.683 [.51,.88] & 0.681 [.51,.88] \\
top\_sv\_excess & 0.609 [.51,.79] & 0.544 [.50,.76] & 0.736$^*$ [.53,.92] \\
norm\_per\_token & 0.598 [.51,.77] & 0.598 [.51,.77] & 0.673 [.52,.85] \\
\bottomrule
\multicolumn{4}{l}{\footnotesize $^*$Rises under the transform; see text.}
\end{tabular}
\end{table}

The full transform raises the spectral gap of every reply, from a hedged
mean of 24.7 to 84.6 or more, and multiplies norm per token by a nearly
constant factor (2.34 to 2.37) in every trial. The spectral gap, the best
feature without the cloak, falls to 0.614 ($p = 0.33$; 0.593 after
removing reply length). Signal rank is 0.766 ($p = 0.020$; 0.755 without
the cloak) and top singular-value excess 0.736 ($p = 0.038$; 0.609
without). But signal rank mostly tracks reply length ($R^2 = 0.77$
against output token count, with or without the cloak), and any cloak
preserves length, because the cache has one row per token. After removing
length, the best feature under the transform is signal rank at 0.705 (top
singular-value excess 0.695), and adjusted for choosing the best of five,
$p = 0.19$; on raw features the best is 0.766, with $p = 0.054$. With
seven confabulations, the best of five raw features reaches 0.77 in 5\%
of label shuffles. The transform removes the one feature that detected
well. Whether a confabulation signal survives it, these data cannot say.

\paragraph{The two $\sigma_1$-dependent features.}
The spectral gap falls (0.903 to 0.614) and top singular-value excess
rises (0.609 to 0.736). Both are consistent with the transform raising
$\sigma_1$, for example through the additive mask; with no ablation and no
recorded parameters, which component does it is not known. With seven
confabulations we cannot tell whether 0.736 is class information
surviving the transform or noise (unadjusted $p = 0.038$; 0.695 after
removing reply length).

% ===================================================================
\section{Discussion}
% ===================================================================

\subsection{What This Means for Dual Use}

Version~1.0 argued that KV-Cloak mitigates the Lyra Technique's dual-use
risk, so that cache-geometry methods could be published while deployed
models stayed protected. The results do not support that. A
rotation-based cloak leaves per-head spectra unchanged, and in the lab's
own layout it shifts and reorders the features but leaves a detectable
signal on both models whose feature is recorded, clearly on Qwen2.5-7B;
it stops only a detector whose threshold is never refitted. Our
implementation of the full transform removed the best feature, but with
seven confabulations we cannot tell whether it removed the signal.
Nothing shown here mitigates the dual-use risk of publishing
cache-geometry methods. A defense would have to be shown on its own
terms: against adaptive adversaries who choose their features after
seeing the cloak, with enough confabulations to measure what remains.

\subsection{A Defender Inside the TEE (an Argument)}

A defender operating inside the TEE could run the Oracle Loop before
obfuscation, so its detection would be the uncloaked detection (AUROC~0.903
on the distilled model) and its steering vectors would be injected in the
model's own basis. That is an argument, not a tested system, and it has a
gap: in KV-Cloak the feature matrix $M$ is fused into the model weights,
so the cache the defender sees is already transformed by $M$, and the
Oracle Loop paper notes that detection under fused $M$ has not been
verified~\citep{lyra2026oracle}. It gives defenders an advantage only if
attackers lose the signal, which the results above do not show.

\subsection{What a Test of the Architecture Would Need}

The deployment architecture version~1.0 recommended remains a proposal,
untested here:

\begin{enumerate}
\item KV-Cloak obfuscation active on all external cache access.
\item Oracle Loop (if deployed) operating inside the TEE before the cache
is obfuscated.
\item Obfuscation key material and Oracle parameters both
TEE-protected.
\item Per-session or per-query key rotation to prevent statistical
recovery from accumulated observations.
\end{enumerate}

The replication used one rotation for every trial, so per-query key
rotation was not exercised there, and the run behind
Table~\ref{tab:real_vs_sim} does not record whether its key changed per
trial. Nothing here measured what it costs an adversary to recover the
signal.

\subsection{Limitations}

\begin{enumerate}
\item \textbf{Small sample size and unaudited labels.} The distilled
model yields only $n=7$ confabulations from 100 prompts. The headline
AUROC 0.903 has bootstrap CI [0.810, 0.980] and survives selection, but
the comparisons under the full transform (Section~\ref{sec:real}) are not
powered, and the labels are a classifier's residual category on replies
cut at 400 tokens and stored as previews (Section~\ref{sec:uncloaked}).
Confirmation on a higher-confab-rate model, with full replies stored, is
needed.

\item \textbf{Feature code not recovered.} The script behind
Tables~\ref{tab:cross_arch} and~\ref{tab:real_vs_sim} has not been
recovered. Its per-trial outputs are in \texttt{data/}; that it used the
lab's row-stacked matrix layout is inferred from its feature names and
from how the features respond to the two rotations. Which layers it
averaged, the full transform's parameters, and what the magnitude-attack
arm did are not recorded.

\item \textbf{Cache versus text untested.} Version~1.0's text-only
baseline cannot be traced or recomputed (Correction note), so nothing here
shows what the cache adds to surface text features.

\item \textbf{TEE trust assumption.} The defensive architecture assumes
TEE integrity. Documented side-channel attacks against SGX
(Van~Bulck et al., 2018; Schwarz et al., 2019) and SEV
(Li et al., 2021) mean this is a hardware trust assumption, not a
cryptographic guarantee. The security reduction is to TEE integrity,
not to a standard hardness assumption.

\item \textbf{Injection test underpowered.} Direct steering, the
positive control, corrected only 0--1 of 10 replies per model, too few to
establish a working baseline. Through-cloak injection results are
therefore uninterpretable: we cannot distinguish ``cloak defeats
injection'' from ``injection was never effective.'' The 30-pair
positive-control run shipped with this paper
(\nolinkurl{data/kv_cloak_positive_control.json}) stored replies without
outcome labels and is not analysed here.

\item \textbf{Cross-architecture significance.} No individual model in
the cross-architecture sweep survives Bonferroni correction for the three
comparisons that have an AUROC. Uncloaked detection on Qwen-14B (0.784)
and Mistral (0.707) is nominal only and should not be read as
cross-architecture replication.

\item \textbf{Adaptive adversary untested.} All tests assume a static,
non-learning attacker. An adversary who can submit probe queries and
observe behavioral outputs could iteratively estimate the rotation
basis. Formal query-complexity bounds are not provided. The threat model
is bounded to single-query, non-adaptive adversaries.
\end{enumerate}

% ===================================================================
\section{Conclusion}
% ===================================================================

Per-head rotation, the simulation of KV-Cloak this paper relied on, does
not remove the signal: on the two models whose feature is recorded, the
feature that detected best without it still detects with it, clearly on
Qwen2.5-7B and marginally on the distilled model, and an adversary who
computes spectra per head sees no change at all. Our implementation of
the full transform removed the best detection feature on the distilled
model, and after removing reply length no other feature detects
significantly; with seven confabulations we cannot say whether a
detectable signal remains. Version~1.0 concluded that KV-Cloak and the
Oracle Loop are complementary, with defenders keeping what attackers lose;
nothing here shows that attackers lose the signal. Whether any reversible
obfuscation blinds cache-geometry detection is open, and a test of it
needs adaptive adversaries, per-query keys, recorded parameters and many
more confabulations.

\section*{Data, Code and Review History}

{\raggedright
Tables~\ref{tab:cross_arch} and~\ref{tab:real_vs_sim}:
\nolinkurl{data/kv_cloak_v2/}, \nolinkurl{data/kv_cloak_real/},
\nolinkurl{data/kv_cloak_cross_arch/} and \nolinkurl{data/kv_cloak_base/},
added in version~1.1, copied unchanged from the lab's KV-Experiments
repository (\nolinkurl{results/oracle_harness/}, commit 8429260, May 2026;
that repository is private). The scripts that wrote them were not
recovered. Replication: \nolinkurl{code/kv_cloak_replication.py} with
\nolinkurl{data/kv_cloak_replication_100.json} and
\nolinkurl{data/kv_cloak_replication_analysis.json}. Positive control:
\nolinkurl{data/kv_cloak_positive_control.json}. Every result in this
version is recomputed from the per-trial files, or read from the run
files where only aggregates exist (Table~\ref{tab:cross_arch}'s
cross-architecture rows, all confidence intervals), by
\nolinkurl{check_correction_claims.py}, in this paper's folder of
\url{https://github.com/Liberation-Labs-THCoalition/published-research};
Methods details carried over from version~1.0 (model architectures, the
classifier's pattern count) are not.
Review history: the lab's remediation plan (23 June 2026; in its
human-review repository,
\nolinkurl{archive/COALITION_REMEDIATION_PLAN_20260623.md}) recommended
withdrawing the rotation claim; the Agni style gate (4 September 2026)
rejected version~1.0
(\nolinkurl{tools/agni/style/kv-cloak-defense-paper.agni_style.json});
Kavi's verification flags (14 September 2026; kept in the lab's internal
notes, not in this repository); for version~1.1, a scoping note
(published-research pull request 10, 20 October 2026), the correction gate
above, and two rounds of the Agni paper and style modes on its drafts
(records in \nolinkurl{review/}).\par}

\bibliographystyle{plainnat}
\section*{Acknowledgments}
We thank Dwayne Wilkes (Sentient Futures / Liberation Labs) for statistical auditing and red-team review.

\bibliography{references}

% ===================================================================
% Appendix: correction record
% ===================================================================
\appendix
\sloppy

\section{Correction Record (October 2026)}
\label{app:record}

This appendix keeps the withdrawn wording of version~1.0 visible. Each
item quotes it verbatim and says where the evidence is; the last items
list other wording that changed.

\begin{enumerate}
\item \textbf{Title.} ``Cache Geometry Under Obfuscation: KV-Cloak as
Defense Against Adversarial KV-Cache Steering.'' It asserted a working
defense (Sections~\ref{sec:shift} and~\ref{sec:real}).

\item \textbf{Abstract, opening.} ``A single invertible linear transform
of the KV cache degrades SVD-based cognitive-state detection from
near-ceiling to near-chance ($p = 5.5 \times 10^{-16}$), while a real
KV-Cloak implementation confirms the defense is conservative (detection
AUROC drops from 0.903 to 0.614).'' The $p$-value is the Wilcoxon test of
Section~\ref{sec:shift}, a shift in one feature's values; under the
rotation the same feature still detects at 0.787
(Section~\ref{sec:shift}); under the full transform, seven confabulations
cannot settle whether the signal is gone (Section~\ref{sec:real}).

\item \textbf{Abstract, finding 1 (text baseline).} ``A text-only
surface-feature baseline achieves 0.755; cache geometry outperforms text
in 96.7\% of bootstrap samples (delta 0.148, 95\% CI [$-$0.008,
0.322])---suggestive but not significant at the 95\% level.'' Results:
``The best text-only feature (average word length) achieves 0.755 [0.560,
0.924]. Cache geometry outperforms text in 96.7\% of bootstrap samples,
but the 95\% CI on the difference includes zero (delta: 0.148 [$-$0.008,
0.322]).'' No run file holds these numbers, and the stored trials keep
only the first 150 characters of each reply.

\item \textbf{Abstract, finding 2.} ``\textbf{Obfuscation defeats
transfer attacks.} [\ldots] A detector calibrated on uncloaked caches
cannot operate on cloaked caches. This is a domain-shift effect: rotation
preserves the underlying geometric structure but changes the coordinate
system, preventing transfer from uncloaked calibration data.'' True only
of a fixed threshold; the ranking survives the rotation, and in a
per-head layout nothing moves (Sections~\ref{sec:replication}
and~\ref{sec:shift}).

\item \textbf{Abstract, finding 3, and the academic limitation on the
magnitude attack.} ``0--2 adverse events per model in 10 trials,
inconsistent with destabilization rates above 30\% (power 0.97 to detect)
but unable to rule out rates up to 20\% at this sample size'' and ``(3
calibration pairs; direct injection produced 0--1/10 corrections)''.
Academic editions: ``At $n=10$ per model, we can detect destabilization
rates above 30\% (power 0.97) but cannot rule out rates up to 20\%.'' Two
events in 10 allow rates up to 51\%; 0.97 is $1 - 0.7^{10}$, the chance of
at least one event; no record gives 3 calibration pairs.

\item \textbf{Abstract, closing.} ``We conclude that KV-Cloak and the
Oracle Loop are complementary: the Oracle Loop provides misalignment
detection and correction; KV-Cloak prevents external adversaries from
weaponizing the same geometric methods. Together, they address the
dual-use risk inherent in publishing cache-geometry techniques.'' And:
``The rotation-simulation result (0.786) represents transfer-attack
failure (domain shift under coordinate change), not signal
degradation.'' The 0.786 (0.787) was measured with both classes under the
same rotation, not by transferring a detector.

\item \textbf{Introduction.} ``We show that it does: obfuscation renders
geometric detection and vector injection incoherent to external
adversaries, while a defender operating inside the trusted execution
environment retains full access to the uncloaked cache and its geometric
properties.''

\item \textbf{Threat Model.} ``[\ldots] but our primary result---that
KV-Cloak blinds external geometric methods---holds regardless of whether
a defensive Oracle is deployed.''

\item \textbf{Method.} ``This preserves within-head singular values while
scrambling cross-head correlations.'' Concatenated along the feature
axis, the heads' singular values are preserved too; stacked along the
token axis, as the lab's extractor stacks them, they are not
(Section~\ref{sec:method_rot}). Also: ``Features are FWL-residualized [\ldots] against output token
count and evaluated via Mann-Whitney AUROC''; the AUROCs reported were computed on raw features.

\item \textbf{Results, feature-space shift.} ``The cloaking is perfectly
detectable by paired comparison---but this is the desired property. An
external adversary's detection model, calibrated on uncloaked caches,
encounters a completely alien feature distribution. Any detection
thresholds, classifier boundaries, or anomaly scores calibrated on
uncloaked data are meaningless on cloaked data.'' (Section~\ref{sec:shift}.)

\item \textbf{Results, injection.} The heading ``Adversarial Injection Is
Incoherent'' and ``The injection becomes noise rather than a targeted
perturbation.'' Untested: the positive control did not work.

\item \textbf{Results, defender.} The heading ``Defensive Oracle Retains
Full Capability'' and ``Detection operates at full capability (AUROC~0.903
on the distilled model).'' An argument, now in the Discussion with the
fused-$M$ gap.

\item \textbf{Results, within-cloak detection.} ``When both confabulating
and hedged responses are cloaked with the same rotation, residual
detection remains possible (AUROC~0.786 on the distilled model, 0.516 on
Mistral---near chance). On Mistral, the cloak effectively eliminates
detection entirely. On the distilled model, residual signal reflects
within-head spectral properties that survive rotation. This
scenario---where both classes share the same obfuscation---is relevant
only to an adversary who controls the cloak key, which the TEE
architecture prevents.'' Within-head spectra survive the rotation
unchanged, but the features here are row-stacked, not within-head
(Section~\ref{sec:method_rot}).

\item \textbf{Results, base model.} ``Detection is weak and heavily
confounded with token count.'' Removing reply length raises it (0.661 to
0.687).

\item \textbf{Results, cross-architecture.} ``Obfuscation reduces
detection on Mistral to chance (0.516) and partially degrades on the
distilled model (0.786 under simulation). Magnitude-based attacks
(direction-agnostic perturbation) produce at most 2 adverse events in 10
trials across all architectures---models are robust to random
perturbation at steering-vector magnitude.''

\item \textbf{Results, simulation vs implementation.} ``Both degrade a
transfer attack---a detector calibrated on uncloaked caches applied to
cloaked caches---rather than destroying the underlying geometric
signal.'' On top singular-value excess: ``This is an artifact of the mask
magnitude, not a genuine detection signal, and would not generalize to
unseen mask instances under per-query key rotation''; the table note
``Anomalous increase; likely artifact of additive mask $A$ inflating top
singular values differentially across behavioral classes''; and ``the
defense works \emph{by} perturbing $\sigma_1$''. Per-query rotation was
not tested, and the transform's parameters were not recorded.

\item \textbf{Discussion.} ``The Lyra Technique's dual-use risk---that
published cache-geometry methods could enable adversarial steering---is
mitigated by KV-Cloak. Cache obfuscation renders both geometric detection
and vector injection incoherent to external adversaries. A model operator
deploying KV-Cloak can publish Oracle Loop methods transparently while
protecting deployed models from exploitation via those same methods.''
``This addresses the responsible-disclosure tension identified by
Watson~(2026): the defense can be published alongside the technique,
enabling the research community to build on the methods while preventing
weaponization against production systems.'' ``Without a defensive Oracle
Loop, KV-Cloak alone prevents external adversaries from deploying
malicious geometric steering.''

\item \textbf{Conclusion.} ``KV-Cloak and the Oracle Loop are
complementary technologies. The Oracle Loop enables detection and
correction of misalignment through KV-cache geometry. KV-Cloak prevents
external adversaries from exploiting those same geometric methods.
Together, they address the dual-use risk inherent in publishing
cache-geometry research: the technique is transparent, the defense is
available, and the architecture ensures that defenders retain
capabilities that attackers lose.''

\item \textbf{Academic editions only.} The abstract of the academic file
\nolinkurl{academic_main.tex}: ``degrades the strongest detection feature
(spectral gap) from 0.903 to 0.614---near chance---confirming that the
defense is stronger than our initial orthogonal-rotation simulation
(0.786) suggested.'' Its Results: ``confirming that per-head orthogonal
rotation alone \emph{understates} the defense. [\ldots] The simulation was
conservative: the actual defense is stronger than we initially
reported.'' A limitation in both academic files: ``The real mechanism
degrades spectral gap detection more than the simulation (0.614 vs
0.786), confirming that the simulation was conservative.'' Another, in
both: ``Until resolved, the scope of the obfuscation-degrades-detection
claim excludes this model'' (Qwen2.5-14B); the claim it scoped is
withdrawn, and the rise is reported in Section~\ref{sec:withincloak}.

\item \textbf{Other changed wording.} Abstract: ``SVD of the key-value
(KV) cache reveals geometric signatures of cognitive states including
confabulation, deception, and sycophancy'', dropped (the deception result
is withdrawn; no sycophancy result is used here). Background:
``features derived from SVD of the key tensor'', now keys and values.
Method: ``we extract five Marchenko-Pastur-derived features from the full
KV cache (encoding + generation phases)'', now five spectral features
averaged over 4 or 7 layers, three of them using the Marchenko--Pastur
edge. Results: ``Signal rank also detects (0.755 [0.570, 0.914])'', now
0.755 [0.578, 0.918] and not significant after correction. Limitations:
``excluding chance but wide'', removed (the intervals are folded and do
not test against chance). Abstract and Limitations: ``no model achieves
significance after Bonferroni correction for four comparisons'', now three
(only three models have an AUROC); the conclusion is unchanged.
\end{enumerate}

\end{document}
